\section{Problem formulation}\label{sec:problem}

\subsection{Model and contract}

Under the pricing measure the asset price $S$ and its variance $v$ follow the \citet{Heston1993} dynamics
\begin{equation}\label{eq:heston}
\dd S_t = r S_t\,\dd t + \sqrt{v_t}\,S_t\,\dd W^{(1)}_t,\qquad
\dd v_t = \kappa(\theta - v_t)\,\dd t + \xi\sqrt{v_t}\,\dd W^{(2)}_t,\qquad
\dd\langle W^{(1)},W^{(2)}\rangle_t = \rho\,\dd t,
\end{equation}
with a constant short rate $r$ and no dividends. (A dividend yield only changes the drift of $S$ and can
be added without difficulty.) Let
\begin{equation*}
A_t = \int_0^t S_u\,\dd u
\end{equation*}
be the running integral of the price. The contract is the fixed-strike arithmetic Asian call with
continuous monitoring over $[0,T]$,
\begin{equation}\label{eq:price}
C(K) = e^{-rT}\,\E\Bigl[\Bigl(\tfrac{1}{T}A_T - K\Bigr)^{+}\Bigr].
\end{equation}

Unless stated otherwise we use the parameters of \citet[Table~2]{Corsaro2019}: $\kappa=1$, $\theta=0.09$,
$\xi=1$, $\rho=-0.3$, $v_0=0.09$, together with $S_0=100$, $r=0.04$ and $T=1$. For these values
$2\kappa\theta/\xi^2 = 0.18$, so the Feller condition fails by a wide margin. The stationary density of the
variance behaves like $v^{2\kappa\theta/\xi^2-1}=v^{-0.82}$ near zero: it is integrable, but it puts a lot
of mass close to the boundary, and the boundary is reached. This is the regime in which we want the solver
to be reliable, and in Section~\ref{sec:regimes} we move $\xi$, $\rho$ and $T$ around it.

\subsection{Similarity reduction}

A direct PDE formulation of \eqref{eq:price} has three space variables, $(S,A,v)$. The price is,
however, homogeneous of degree one in $(S,A,K)$, and following \citet{RogersShi1995} and
\citet{Vecer2001,Vecer2002} we write the time-$t$ value as
\begin{equation}\label{eq:ansatz}
V(t,S,A,v) = S\,w(t,x,v),\qquad x = \frac{K - A/T}{S}.
\end{equation}
Substituting \eqref{eq:ansatz} into the pricing equation for $V$, the terms $rV$ and $rS\,w$ cancel and
one is left with an equation for $w$ in two space variables,
\begin{equation}\label{eq:pde}
\frac{\partial w}{\partial t}
+ \tfrac12 v x^2 \frac{\partial^2 w}{\partial x^2}
- \rho\xi v x \frac{\partial^2 w}{\partial x\,\partial v}
+ \tfrac12 \xi^2 v \frac{\partial^2 w}{\partial v^2}
- \Bigl(r x + \frac1T\Bigr)\frac{\partial w}{\partial x}
+ \bigl(\kappa(\theta - v) + \rho\xi v\bigr)\frac{\partial w}{\partial v} = 0,
\end{equation}
with terminal condition
\begin{equation}\label{eq:terminal}
w(T,x,v) = \max(-x,0),
\end{equation}
since $(A_T/T-K)^+ = S_T\,(-x_T)^+$. At $t=0$ we have $A_0=0$, hence $x_0=K/S_0$ and
\begin{equation}\label{eq:readout}
C(K) = S_0\,w\bigl(0,\,K/S_0,\,v_0\bigr).
\end{equation}
\citet{EwaldMenkensTing2013} arrive at an equivalent equation through a change of measure to the stock
as num\'eraire; in that picture $w$ is a price in units of the stock, which explains why no discounting
term remains. Three properties of \eqref{eq:pde} shape everything that follows.

The first is that its coefficients do not depend on $t$. The discretised operator is therefore the same
at every time step, and any matrix factorisation can be reused for the whole time march.

The second is that the strike does not appear in \eqref{eq:pde}--\eqref{eq:terminal} at all. It enters
only through the point $x_0=K/S_0$ at which the solution is read in \eqref{eq:readout}, so one solution
prices all strikes whose $x_0$ lies inside the computational domain.

The third is that the equation contains the one-dimensional problem as a special case. If $\xi\to0$ and
$v_0=\theta=\sigma^2$, the variance stays at $\sigma^2$ and \eqref{eq:pde} becomes the Rogers--Shi
equation $w_t + \tfrac12\sigma^2x^2w_{xx} - (rx + 1/T)w_x = 0$ for geometric Brownian motion.

\subsection{Degeneracies and boundaries}\label{sec:boundaries}

\paragraph{The variance boundary $v=0$.} Every second-order term in \eqref{eq:pde} carries the factor $v$.
On $v=0$ the equation therefore reduces to the first-order equation
\begin{equation}\label{eq:v0}
\frac{\partial w}{\partial t} - \Bigl(rx+\frac1T\Bigr)\frac{\partial w}{\partial x} + \kappa\theta\,\frac{\partial w}{\partial v} = 0 .
\end{equation}
Since $\kappa\theta>0$, the drift of the variance points into the domain, and by the theory of equations
with nonnegative characteristic form \citep{OleinikRadkevich1973} no boundary condition may be imposed
on $v=0$; the value there is determined by the equation. For the Heston model this is the situation
analysed by \citet{EkstromTysk2010}. The practical consequence for a numerical method is that the
boundary line must carry the degenerate equation \eqref{eq:v0} itself, and not an artificial condition.
This holds whether or not the Feller condition is satisfied; what changes when it fails is how much
probability mass sits near the boundary.

\paragraph{The line $x=0$.} The diffusion coefficient $\tfrac12vx^2$ also vanishes on $x=0$, which is an
interior line of the domain, and the terminal condition \eqref{eq:terminal} has a kink there. The kink is
closer to the price than it looks. Along the drift alone, $x$ moves according to
$\dd x = -(rx+1/T)\,\dd t$, so that $x_T = (x_0 + 1/(rT))e^{-rT} - 1/(rT)$. For the at-the-money readout
$x_0=1$ and our parameters this gives $x_T\approx-0.02$: the characteristic through the readout point ends
almost exactly on the kink, in a region where the diffusion in $x$ is small.

\paragraph{Far field.} For $x\to-\infty$ the option is deep in the money and $w$ becomes linear in $x$;
for $x\to+\infty$ it is worthless. We truncate to $x\in[x_{\min},x_{\max}]$ and impose
\begin{equation}\label{eq:bcx}
\frac{\partial^2 w}{\partial x^2}(t,x_{\min},v) = 0,\qquad w(t,x_{\max},v) = 0 .
\end{equation}
For large variance we truncate at $v_{\max}$ and impose
\begin{equation}\label{eq:bcv}
\frac{\partial^2 w}{\partial v^2}(t,x,v_{\max}) = 0 .
\end{equation}
This condition allows $w$ to grow linearly in $v$ near the cut-off. In our tests it was much less sensitive
to the position of $v_{\max}$ than the more common $\partial w/\partial v=0$: at the base parameters,
moving $v_{\max}$ from $1.5$ to $4$ changed the at-the-money price by $7.0\times10^{-4}$ under the
Neumann condition and by $1.4\times10^{-6}$ under \eqref{eq:bcv} (at most $8\times10^{-6}$ over
$K=90,100,110$). The default domain is $x\in[-1,5]$,
$v\in[0,1.5]$. Section~\ref{sec:regimes} shows that this default has to be widened for some parameters.

\subsection{A geometric twin with an exact price}\label{sec:twin}

To test the numerical machinery against an exact answer we also solve the geometric-average version of
the contract,
\begin{equation*}
C_G(K) = e^{-rT}\,\E\bigl[(G_T - K)^+\bigr],\qquad G_T = \exp\Bigl(\frac1T\int_0^T \ln S_u\,\dd u\Bigr).
\end{equation*}
With $a(t) = (T-t)/T$ the variable
\begin{equation*}
\zeta_t = \frac1T\int_0^t \ln S_u\,\dd u + a(t)\ln S_t
\end{equation*}
satisfies $\dd\zeta_t = a(t)\bigl[(r - v_t/2)\,\dd t + \sqrt{v_t}\,\dd W^{(1)}_t\bigr]$; the term in
$\ln S_t$ cancels because $a'(t)=-1/T$. The price $u(t,\zeta,v)$ therefore solves
\begin{equation}\label{eq:pdegeo}
u_t + a(r - \tfrac v2)u_\zeta + \kappa(\theta-v)u_v + \tfrac12 a^2 v\,u_{\zeta\zeta} + \rho\xi a v\,u_{\zeta v}
+ \tfrac12\xi^2 v\,u_{vv} - r u = 0,\qquad u(T,\zeta,v) = (e^\zeta-K)^+,
\end{equation}
and $C_G(K) = u(0,\ln S_0,v_0)$. The variance part of \eqref{eq:pdegeo} is the same as in \eqref{eq:pde},
including the degeneracy at $v=0$, so the same boundary analysis applies. Two differences matter. The
coefficients now depend on time through $a(t)$, which we will see costs a factorisation per time step.
And the kink of the payoff lies at $\zeta=\ln K$, which for $K=S_0$ coincides with the readout point;
moreover the accumulated variance $\int_t^T a^2 v\,\dd s \approx v(T-t)^3/(3T^2)$ is cubically small near
maturity, so the kink is smoothed slowly. The geometric twin is thus a hard test of the time and space
discretisation, but it does not test the arithmetic reduction \eqref{eq:ansatz}. On the far boundaries we
use $u=0$ at $\zeta_{\min}$, the condition $u_{\zeta\zeta}-u_\zeta=0$ at $\zeta_{\max}$ (exact for any
function $Ae^\zeta+B$), and $u_v=0$ at $v_{\max}=4$.
